\section{Experimental Protocol}
\label{sec:protocol}

\subsection{Data, fitting and evaluation exposure}
HRSID~\cite{hrsid} supplies 2,937 official training images for the detectors. We cross-fit the fractional occupancy head in five folds on FIT109, grouping tiles by source filename prefix, and refit it on the source data before transfer. DEV705 contains 705 images and 1,782 hulls from 28 filename-prefix groups. We found no exact encoded, decoded, or prefix overlap with FIT. We use DEV705 to compare endpoint writers during development. The HRSID official test allocation has 1,962 images and 5,922 annotations and was previously used for model selection. The historical evaluation includes the full allocation, with separate sensitivity checks for four empty rasterized annotations and three image identities involved in exact overlaps.

Instance-RSDD~\cite{rsdd} provides independent author annotations for the frozen HRSID fits. The obtained test allocation contains 740 chips, 872 hulls and 112 parent acquisitions. We prespecified a Gaofen-3 subset with 583 chips, 631 hulls and 76 parent acquisitions because Gaofen-3 is absent from detector fitting. The authors' own training and test allocations share parents. The author training allocation supplies a separate 74-chip reserve with 81 hulls from 13 parents absent from the 740-chip allocation. We froze the source checkpoint and signed comparison before its first prediction and score. A decoded-image check found zero exact duplicates against HRSID fitting and the earlier RSDD test. The nested 42-chip view excludes shared sensor--date keys. A joint screen for shared sensor--date keys or same-sensor rounded centers leaves 17 eligible matches from five parents, below the original coverage criterion. The later boundary analyses revisit scored images. Geographic and orbital separation remain unverified.

PSeg-SSDD~\cite{ssdd} and the corresponding SPSD masks assign separate instance boundaries to the same 232 historical SSDD images. The respective annotations include 546 and 547 hulls. The 222-image sensitivity removes ten exact decoded overlaps with the PSeg authors' training allocation. For the two-reference experiment, 732 SSDD training images form the head-fitting partition and 183 source-development images form the evaluation partition. The latter include 180 exact decoded-image groups, with 411 PSeg and 410 SPSD hulls. Decoded groups are disjoint between fitting and evaluation. This evaluation partition informed method development. During transfer, SAR-SSP~\cite{slgsea,airsarship} supplies 139 raster ship instances in 232 validation chips derived from 31 Gaofen-3 AIR-SARShip-1.0 scenes. The SAR-SSP authors' training and validation chips share parent acquisitions. Our transferred detectors and heads are frozen from other datasets. SAR-SSP has prior evaluation exposure in the project. Prospective scoring on the separately registered Umbra acquisitions awaits independent reference masks.

\subsection{Fixed detectors and controls}
D-FINE-seg-S~\cite{dfine} starts from the official COCO weights and uses the 24th-epoch exponential moving average after 17,640 HRSID optimizer updates. Its native inference uses 640-pixel inputs, 300 queries, box nonmaximal suppression at IoU 0.70 and mask threshold 0.50. The two RF-DETR~\cite{rfdetr} fits use the same COCO initialization, HRSID training images and update count, with separate detector seeds 5202 and 5203. RF-DETR v1.12.0 uses the author 576-pixel validation transform and 164 queries, with joint top-100 selection before retaining ship-class slots. Each comparison holds the detector's candidate stream and score order fixed. The fixed signed-count comparison also holds the addition and removal counts and writable endpoints fixed while changing the positions written. Section~\ref{sec:variable-sign} separately varies the division of a fixed total count between additions and removals.

The scalar occupancy head has 1,321 trainable parameters and uses 200 AdamW updates with source-group sampling (Section~\ref{sec:occupancy}). The dual-reference head has 1,346 parameters and uses 1,000 updates with 64 candidates per update. The PointRend adaptation used for reserve timing has 198,402 parameters and uses 200 updates. Frozen detector inference runs on CUDA. The RF-DETR verification runtime uses PyTorch 2.10.0 and CUDA 13.0. Local writing and PointRend timing use one numerical thread on an Intel Core i7-13700KF with inputs in memory. Phoenix XL1 uses four two-thread CPU workers on DGX Spark, batches of up to two masks and five author iterations. These are implementation settings for the observed comparisons.

\subsection{Detection scores and fixed-match geometry}
\label{sec:metrics}
COCO scoring uses the original detection scores and at most 100 detections per image. At raster mask IoU threshold \(t\), we write \(P_t(j)\) and \(R_t(j)\) for cumulative precision and recall through ranked result \(j\). For eligible references \(N_G\), these are \(\mathrm{TP}_t(j)/[\mathrm{TP}_t(j)+\mathrm{FP}_t(j)]\) and \(\mathrm{TP}_t(j)/N_G\), respectively. With \(\mathcal T=\{0.50,0.55,\ldots,0.95\}\), \(\mathcal R=\{0,0.01,\ldots,1\}\) and the interpolated precision \(P_t^*(r)=\max_{j:R_t(j)\geq r}P_t(j)\), we compute
\begin{equation}
\begin{aligned}
 \mathrm{AP}&=\frac{\sum_{t\in\mathcal T}\sum_{r\in\mathcal R}P_t^*(r)}
 {|\mathcal T||\mathcal R|},\\
 \mathrm{AR}_{100}&=\frac{1}{|\mathcal T|}\sum_{t\in\mathcal T}\max_j R_t(j).
\end{aligned}
\label{eq:coco-metrics}
\end{equation}
We set \(P_t^*(r)=0\) at unreached recall levels. AP75 and AP90 use their respective IoU thresholds and small-mask AP uses the COCO native area bin through \(32^2\). COCO assigns matches separately for every writer and threshold.

For geometry we freeze the original nonignored, score-ordered one-to-one assignment at raster IoU 0.50. The primary small-hull analyses use original predicted-mask areas in \([32,1024)\) native pixels, including matches with zero edits. Given edited mask \(M'\) of area \(P'\) and polygon \(G\) of area \(T\), we intersect the union of native unit cells with \(G\) to obtain area \(I'\). The resulting measures are
\begin{equation}
 \begin{aligned}
  \mathrm{IoU}_{\mathrm{cont}} &= \frac{I'}{P'+T-I'},\\
  F_{\mathrm{P}} &= \frac{P'-I'}{T},\qquad F_{\mathrm{N}}=\frac{T-I'}{T}.
 \end{aligned}
 \label{eq:fixed-geometry-metrics}
\end{equation}
For raster references, IoU uses foreground cell counts. Boundary IoU~\cite{boundaryiou} compares the inward rings \(B_d(X)=X\setminus\mathrm{erode}_{3\times3}^{\,d}(X)\), using zero exterior padding. We compute the IoU of \(B_d(M')\) and \(B_d(G_{\rm r})\) with \(d=1\) as the primary width and \(d=2\) as a sensitivity. The non-saturated boundary subset requires interior pixels in both the native mask and the reference after erosion.

For paired hull differences \(\Delta_{gi}\) on \(G_b\) contributing groups, with \(n_g\) hulls in group \(g\), we estimate
\begin{equation}
 \overline{\Delta}
 =\frac{1}{G_b}\sum_{g=1}^{G_b}\frac{1}{n_g}
  \sum_{i=1}^{n_g}\Delta_{gi}.
 \label{eq:paired-acquisition-effect}
\end{equation}
We average head effects before group resampling when heads share the native assignments. Reported 95\% intervals use 4,096 paired group resamples for the main contrasts. Reserve sensitivity resamples sensor--date blocks while retaining equal parent weights. For the dual-reference evaluation, all-hull raster IoU assigns zero to native misses and averages hulls within image before averaging images within decoded groups. The SAR-SSP version includes all labeled hulls, with zero change for native misses, and gives equal weight to the 23 ship-bearing parent acquisitions.

\section{Quantitative Results}
\label{sec:quantitative}

\subsection{Endpoint selection on HRSID}
\begin{table}[t]
 \caption{D-FINE E24 on HRSID DEV705. AP and recall use all 70,500 candidates. Continuous fixed-match IoU weights 27 source-prefix groups equally across 1,447 eligible native matches. Scores are in points. Matched counts come from distance pairing. Natural counts come from the independent writer.}
 \label{tab:dev}
 \centering\footnotesize
 \setlength{\tabcolsep}{3pt}
 \begin{tabular}{@{}lrrrr@{}}
 \toprule
 Writer & AP & AP75 & AR100 & IoU\\
 \midrule
 Native mask & 48.160 & 51.507 & 54.018 & 75.798\\
 Distance pairing & 48.916 & 53.453 & 54.792 & 76.258\\
 Independent, matched & 49.726 & 55.285 & 55.617 & 76.787\\
 Native probability, natural & 47.951 & 50.616 & 53.765 & 75.652\\
 Independent, natural & 50.715 & 57.760 & 56.695 & 77.217\\
 \bottomrule
 \end{tabular}
\end{table}

On DEV705, independent selection at the distance-paired counts gains 0.811 AP and 0.529 continuous IoU points (\([0.455,0.597]\)), with 83,059 swaps across 20,741 candidates (Table~\ref{tab:dev}). At its natural counts it makes 194,052 swaps and gains 2.555 AP and 1.419 IoU points over the native masks (\([1.265,1.587]\)). AP90 increases from 0.145 to 0.223 points. Ranking the same natural-count writes by the frozen detector probabilities yields 47.951 AP, so the positions matter even when this edit count is held fixed. The natural-count versus native comparison changes the extent of editing. On the two RF-DETR fits, the matched-count contrast on DEV705 gives \(+0.906\) and \(+0.838\) AP, with corresponding IoU effects \(+0.639\) (\([0.548,0.743]\)) and \(+0.583\) (\([0.488,0.678]\)).

On the previously used HRSID official test allocation, the three-head mean for RF-DETR is 54.433 AP for the native masks, 54.227 for distance pairing and 55.444 for independent selection at matched counts. The latter improves on pairing by 1.216 AP, 1.958 AP75 and 1.379 small-mask AP points. Individual AP effects are \(+1.052,+1.297,+1.299\). The main AP effect ranges from 1.207 to 1.216 under the overlap and empty-raster sensitivities. Natural-count selection gains 1.021 AP over the native masks. This official-split comparison uses archived detector fit 5202. Under the same 24-epoch DEV705 schedule, the native RF-DETR raster fit scores 51.434 AP and the center-supervised fit scores 56.072. These reference values represent the combined effects of the detector, supervision and training recipe.

\subsection{Frozen transfer and boundary geometry}
\begin{table*}[t]
 \caption{Independent selection minus distance pairing at fixed signed counts on Instance-RSDD. Effects are in points. Match and parent counts describe the eligible original assignments and represented parent scenes. Intervals come from paired group resampling. The reserve 5202 comparison uses head 5217.}
 \label{tab:primary-quant}
 \centering\footnotesize
 \setlength{\tabcolsep}{5pt}
 \begin{tabular}{@{}llrrrrr@{}}
 \toprule
 View, detector & Head seeds & Matches/parents & \(\Delta\)AP & \(\Delta\)continuous IoU & \(\Delta\)Boundary IoU (1 px) & \(\Delta\)Boundary IoU (2 px)\\
 \midrule
 Gaofen-3, 5202 & 5217, 5218 & 425/68 & \(+0.483\) & \(+0.259\,[0.194,0.322]\) & \(+0.466\,[0.223,0.702]\) & \(+0.467\,[0.278,0.653]\)\\
 Gaofen-3, 5203 & 5217--5219 & 425/67 & \(+0.299\) & \(+0.196\,[0.136,0.253]\) & \(+0.177\,[-0.025,0.376]\) & \(+0.255\,[0.107,0.404]\)\\
 Reserve, 5202 & 5217 & 47/12 & \(+0.436\) & \(+0.318\,[0.217,0.436]\) & \(+0.701\,[0.180,1.233]\) & \(+0.617\,[0.236,1.001]\)\\
 \bottomrule
 \end{tabular}
\end{table*}

Table~\ref{tab:primary-quant} compares two detector fits on the same Gaofen-3 images. The fixed-count AP effects are \(+0.482,+0.483\) for the first fit and \(+0.240,+0.280,+0.378\) for the second. Both fits improve group-weighted continuous IoU. At one pixel, the boundary interval excludes zero for 5202 and spans zero for 5203. At two pixels both intervals exclude zero. The reserve effect for 5202 is positive at both widths. Its one-pixel sensor--date block interval is \([0.231,1.243]\). All eligible native/reference masks in these comparisons retain interior cells at both widths. At natural counts, the two 5202 heads improve AP over native by 0.587 and 0.517 on Gaofen-3, and AP75 by 1.438 and 1.230. AP90 decreases by 0.094 and 0.107, respectively.

The 74-chip reserve also gives \(+0.306\) AP and \(+0.267\) continuous IoU (\([0.161,0.374]\)) for fit 5203 at matched counts. In the nested 42-chip sensor--date-separated view, the AP effects are \(+0.840\) and \(+0.529\) for the two fits. Corresponding continuous IoU effects are \(+0.364\) (\([0.250,0.478]\)) and \(+0.326\) (\([0.166,0.460]\)). These are repeated fits on common imagery and estimates from related reserve views. On the PSeg official allocation, the same endpoint contrast gives \(+0.081\) AP and \(+0.106\) continuous IoU (\([0.055,0.157]\)). Against SPSD on the same imagery the AP effect is \(-0.007\). On the same SSDD images, the natural-count writer changes AP relative to native by \(+0.344\) against PSeg and \(-0.491\) against SPSD. The endpoint gains therefore depend on the evaluation reference as well as the acquisition cohort.

\subsection{Strong writers at comparable budgets}
PointRend~\cite{pointrend} uses HRSID FIT109 RF-DETR features, 200 updates at 196 targets per update and 3,136 feature evaluations per region of interest. The matched writer gives both methods the same endpoints and signed counts. On the 74-chip Instance-RSDD reserve, this yields 31,391 swaps per arm. The occupancy writer gives \(+0.477\) AP and \(+0.395\) continuous IoU (\([0.204,0.620]\)). The one-pixel Boundary IoU effect is \(+0.829\) (\([-0.056,1.750]\)) and the two-pixel effect is \(+0.801\) (\([0.285,1.370]\)). On Gaofen-3, the occupancy writer leads PointRend by 0.312 AP on the common capped support and 0.390 AP on the common complete boundary support. When allowed to paste its whole region of interest, PointRend leads complete-boundary occupancy by 0.200 AP, with 1.744 points lower AR100. That comparison changes writable support.

A frozen Phoenix XL1~\cite{phoenix} supplies a separate whole-grid control. Holding the posterior and editable candidates fixed, an exact-count whole-grid write raises AP from 49.085 to 58.478 on the reserve and from 13.922 to 16.447 on PSeg-SSDD, relative to the author threshold writer on that grid. The respective counts fall from 2,042,908 to 59,800 and from 3,391,327 to 147,540 flips. In the same order, continuous IoU increases by 9.462 (\([6.329,13.283]\)) and 4.001 points (\([2.644,5.382]\)). At the same narrow support and signed counts, Phoenix attains 58.976 AP on the reserve versus 58.668 for the occupancy head mean, with a continuous-IoU interval spanning zero. These comparisons distinguish the effect of suppressing excess writes from any advantage of a particular ranking model.

\subsection{Variable signed allocation and frozen transfer}
\label{sec:variable-sign-results}

\begin{table}[t]
 \caption{Variable minus balanced signed allocation on the 183-image source-development split. Both writers use the same fitted head, candidates, support and total edit count. Effects are in points.}
 \label{tab:dual-sign-seeds}
 \centering\footnotesize
 \setlength{\tabcolsep}{4pt}
 \begin{tabular}{@{}lrrrr@{}}
 \toprule
 & \multicolumn{2}{c}{PSeg} & \multicolumn{2}{c}{SPSD}\\
 \cmidrule(lr){2-3}\cmidrule(l){4-5}
 Head seed & \(\Delta\)AP & \(\Delta\)IoU & \(\Delta\)AP & \(\Delta\)IoU\\
 \midrule
 5297 & \(+2.461\) & \(+0.993\) & \(+1.782\) & \(+0.580\)\\
 5298 & \(+2.465\) & \(+1.037\) & \(+1.615\) & \(+0.584\)\\
 5299 & \(+1.971\) & \(+0.882\) & \(+1.856\) & \(+0.555\)\\
 Mean & \(+2.299\) & \(+0.971\) & \(+1.751\) & \(+0.573\)\\
 \bottomrule
 \end{tabular}
\end{table}

Each fit evaluates 18,300 original candidate slots and 115,384 flips per writer. With the same estimated-IoU objective and head, widening the menu of admissible signed splits improves both AP and all-hull raster IoU on both references for all three seeds (Table~\ref{tab:dual-sign-seeds}). The group intervals for the mean are \([0.740,1.207]\) for PSeg and \([0.319,0.840]\) for SPSD. The seeds share scenes and labels. For seed 5297, using the mean of the two reference objectives in place of their minimum yields \(+0.498\) AP and \(+0.312\) IoU on PSeg, and \(-0.044\) AP and \(-0.073\) IoU on SPSD. The stable comparison here concerns the signed split, with a trade-off between the two objective aggregations.

\subsection{Frozen signed allocation on SAR-SSP}
\label{sec:sar-ssp}
The frozen seed-5297 writer on SAR-SSP yields \(-0.139\) AP and \(-0.212\) all-hull IoU (\([-0.547,0.104]\)) relative to balanced allocation at the same 147,118 flips. Of 139 labeled hulls, 88 have an original native match. The writers differ at 48 of those matches across 20 groups. False-positive area increases by 1.378 points of reference area (\([0.750,2.067]\)) while missed area decreases by 0.653 points (\([-1.079,-0.328]\)). In the separate HRSID-trained endpoint comparison with fixed signed counts, RF-DETR gives \(-0.005\) AP and \(+0.101\) IoU (\([-0.066,0.335]\)), and D-FINE gives \(+0.076\) AP and \(+0.022\) IoU (\([-0.193,0.238]\)).
